% Generated draft: requires mathematical and source review.
\documentclass[11pt]{article}
\usepackage[margin=25mm]{geometry}
\usepackage{fontspec}
\setmainfont{Latin Modern Roman}
\usepackage{amsmath,amssymb,mathtools,unicode-math}
\setmathfont{Latin Modern Math}
\usepackage{enumitem,needspace,xurl,hyperref}
\setlength{\parindent}{0pt}
\setlength{\parskip}{5pt}
\setlength{\emergencystretch}{3em}
\newcommand{\N}{\mathbb{N}}
\newcommand{\Z}{\mathbb{Z}}
\newcommand{\Q}{\mathbb{Q}}
\newcommand{\R}{\mathbb{R}}
\newcommand{\C}{\mathbb{C}}
\newcommand{\F}{\mathbb{F}}
\newcommand{\A}{\mathbb{A}}
\newcommand{\PP}{\mathbb{P}}
\newcommand{\E}{\mathbb{E}}
\newcommand{\eps}{\varepsilon}
\newcommand{\dd}{\,\mathrm{d}}
\begin{document}
\section*{UnsolvedMath ID: 1344}
\textbf{Source problem code:} catalogue entry
\subsection{Definitions and mathematical statement}
For a group $G$, let $\operatorname{Th}(G)$ denote its complete first-order theory, using the pure group language as the standard interpretation. A complete theory $T$ in a countable language is $\aleph_0$-stable, or $\omega$-stable, if for every countable parameter set $A$ in a model of $T$, the set $S_1(A)$ of complete one-variable types over $A$ is countable. A complete type records a consistent complete choice of formulas in that variable with parameters from $A$, allowing realizations in elementary extensions.

A group is abstractly simple if it is nontrivial and has no normal subgroups other than itself and the identity subgroup. This is stronger than merely having no proper nontrivial \emph{definable} normal subgroup. An algebraic group over a field $k$ is an algebraic variety with regular multiplication and inversion; write $H(k)$ for its group of $k$-rational points. A field $k$ is algebraically closed if every nonconstant polynomial in $k[t]$ has a root in $k$.

\textbf{Imported $\omega$-stable algebraicity question.} For every \emph{infinite} abstractly simple group $G$ such that $\operatorname{Th}(G)$ is $\omega$-stable, does there exist an algebraically closed field $k$ and an algebraic group $H$ over $k$ such that
\[
 G\cong H(k)
\]
 as abstract groups? The infinitude condition comes from the supplied background, although it is omitted from the short statement. The background further describes the algebraic target as simple. Conventions concerning simple algebraic groups and their possible finite centers must be made compatible with the abstract simplicity of $G$; the record does not state a central-quotient convention or require the isomorphism to be definable.

The usual finite-Morley-rank Cherlin--Zilber formulation is a related but distinct statement. Finite Morley rank concerns finiteness of the recursively defined rank of the underlying definable set: a nonempty definable set has rank at least $0$, and it has rank at least $\alpha+1$ if it contains infinitely many pairwise disjoint definable subsets of rank at least $\alpha$, with the usual limit-ordinal clause and with parameters allowed. The source uses $\omega$-stability, not a finite-Morley-rank hypothesis, and this draft does not replace one by the other.

The literature triage dated 2026-08-17 describes the general algebraicity programme as unresolved and explicitly warns against conflating these formulations. It reports algebraicity for simple $K^*$-groups of finite Morley rank of even or mixed type, and says odd/degenerate cases and nonalgebraic bad-group possibilities remain unresolved in the relevant general setting. Those technical subclasses are contextual progress, not hypotheses added to the imported $\omega$-stable question. The supplied background also cites Zilber's work on Zariski geometries as evidence, without supplying a particular reference. The registry's \emph{open} and triage's \emph{partially solved} classifications are retained as supplied dated assessments, not independently verified current claims.
\subsection{Short English statement}
Ask whether every infinite simple group whose first-order theory is $\omega$-stable is isomorphic to the points of an algebraic group over an algebraically closed field. The record uses this stronger stability-based formulation and explicitly distinguishes it from the usual finite-Morley-rank Cherlin--Zilber conjecture.
\subsection{Sources}
\begin{itemize}
\item UnsolvedMath local registry, data/unsolvedmath/problems.json, permanent ID 1344, source code catalogue entry, ``Cherlin-Zilber Conjecture,'' including its background's infinitude condition and literature triage dated 2026-08-17.
\item K. Tent, ``From the Cherlin-Zilber Conjecture via sharply 2-transitive groups to the Burnside problem,'' arXiv:2606.18207 (2026), \url{https://arxiv.org/abs/2606.18207}. The supplied review cites it for the open finite-Morley-rank algebraicity programme.
\item G. Cherlin, ``Simple Groups of Finite Morley Rank,'' Rutgers research page, \url{https://sites.math.rutgers.edu/~cherlin/FMR/}. The supplied review cites even/mixed-type algebraicity and unresolved odd/degenerate cases.
\item O. Deloro and E. Jaligot, ``Superstability and central extensions of algebraic groups,'' \emph{Ann. Pure Appl. Logic} 167 (2016), 75--103, \url{https://doi.org/10.1016/j.apal.2015.08.004}. The supplied review cites the $\omega$-stable formulation and its relationship with the broader algebraicity conjecture.
\item Ulam AI and the UnsolvedMath Contributors, \emph{UnsolvedMath: A Curated Collection of Open Mathematics Problems}. Dataset curation and original metadata are supplied under CC BY 4.0; referenced source material retains its own terms. No external reference was independently consulted for this draft.
\end{itemize}
\end{document}
